\section*{Capítulo \ref{ch:b3:field-lab-methods-3} --- Métodos de campo e de laboratório III}
% Values from the Book 5 scratch calculator calc/ch31.py.

\begin{solution}{exo:b3:field-lab-methods-3:1}
(a)~$\qty{0.01}{\milli\gal} = \qty{e-7}{m.s^{-2}}$, e
$\num{e-7}/\num{9.81} = \num{1.0e-8}$: uma parte em cem milhões.

(b)~$0.01/0.3086 = \qty{0.032}{m}$, cerca de \qty{3.2}{cm}.

(c)~Juntos, os dois efeitos alteram a gravidade em
\qty{0.1967}{\milli\gal.m^{-1}}, a diferença $0.3086 - 0.1119$, logo
a altura deve variar em $0.01/0.1967 = \qty{5.1}{cm}$.
\end{solution}

\begin{solution}{exo:b3:field-lab-methods-3:2}
$\Delta m = \tfrac43\pi(\qty{50}{m})^3\times\qty{500}{kg.m^{-3}} =
\qty{2.62e8}{kg}$. O pico é $G\Delta m/z^2 =
\num{6.674e-11}\times\num{2.62e8}/\num{e4} = \qty{1.75e-6}{m.s^{-2}}$,
isto é \qty{0.175}{\milli\gal}. A meia-largura é $x_{1/2} = z/1.305 =
\qty{77}{m}$.
\end{solution}

\begin{solution}{exo:b3:field-lab-methods-3:3}
(a)~Com $AM = MN = NB = a$, $K = 2\pi/(1/a - 1/2a - 1/2a + 1/a) =
2\pi a = \qty{62.83}{m}$, e $\rho_a = 62.83\times0.239/0.1 =
\qty{150}{\ohm\metre}$. (b)~$K = \pi(15^2 - 5^2)/(2\times5) = \pi\times
200/10 = \qty{62.83}{m}$. Os eletrodos ficam então em $0$, $10$, $20$
e \qty{30}{m}, na ordem AMNB: trata-se do arranjo de Wenner com espaçamento de
\qty{10}{m}.
\end{solution}

\begin{solution}{exo:b3:field-lab-methods-3:4}
(a)~$V = 0.2998/\sqrt5 = \qty{0.134}{m.ns^{-1}}$. (b)~$t =
2\times10/0.134 = \qty{149}{ns}$. (c)~$\delta = 0.0053\sqrt5/\num{e-3} =
\qty{11.9}{m}$, e uma profundidade útil de cerca de \qty{36}{m}. (d)~Na argila
$\delta = 0.0053\sqrt{20}/0.01 = \qty{2.4}{m}$, uma profundidade útil de cerca de
\qty{7}{m}, e uma velocidade de apenas \qty{0.067}{m.ns^{-1}}. A
aproximação ainda vale ali: a \qty{100}{MHz},
$\omega\varepsilon = 2\pi\times\num{e8}\times20\times\num{8.85e-12} =
\qty{0.11}{S.m^{-1}}$, dez vezes $\sigma$.
\end{solution}

\begin{solution}{exo:b3:field-lab-methods-3:5}
(a)~O cilindro tem $\lambda = \pi R^2\Delta\rho$ por unidade de comprimento, e
fora dele, sua atração é a de uma massa linear sobre o seu eixo (o teorema de Gauss
aplicado a um cilindro coaxial): $\Delta g = 2G\pi R^2\Delta\rho\,z/(x^2
+ z^2)$. Metade do pico é atingida quando $x^2 + z^2 = 2z^2$, logo $z =
x_{1/2}$. (b)~A esfera está a $1.305\times100 = \qty{130}{m}$, o
cilindro a \qty{100}{m}: o mesmo perfil, interpretado como esfera, situa o
corpo \qty{30}{\percent} mais profundo. A anomalia de uma esfera decai mais depressa com
a distância, logo uma esfera precisa estar mais profunda para ser tão larga.
\end{solution}

\begin{solution}{exo:b3:field-lab-methods-3:6}
(a)~Cada ponto representa $\qty{100}{m}\times\qty{100}{m} =
\qty{e4}{m^2}$, logo $\iint\Delta g\,\mathrm{d}A \approx
\num{-438e-5}\times\num{e4} = \qty{-43.8}{m^3.s^{-2}}$, e $\Delta m
= -43.8/(2\pi\times\num{6.674e-11}) = \qty{-1.04e11}{kg}$. (b)~$1.04/1.74
= \qty{60}{\percent}$. O restante está nos flancos da anomalia fora
do mapa: o granito atinge uma profundidade de mais de um quilômetro, e a
anomalia de sua parte profunda decai lentamente com a distância. Estender o mapa
o bastante, ou usar um modelo, recupera-a.
\end{solution}

\begin{solution}{exo:b3:field-lab-methods-3:7}
(a)~Para um dique delgado, a anomalia cai à metade do pico em $x = \pm z$:
sua largura total à meia altura do pico é $2z = \qty{60}{m}$, e três
leituras sobre ela exigem espaçamento de no máximo \qty{20}{m}. (b)~$2\times
220/3 = \qty{147}{m}$. (c)~O espaçamento necessário entre as estações depende da
profundidade do corpo mais raso de interesse: um levantamento projetado para o dique
profundo, com espaçamento de \qty{150}{m}, distorceria a anomalia do raso, que
tem apenas \qty{60}{m} de largura; um levantamento projetado para o corpo raso também enxerga o
profundo, com muito mais leituras do que o necessário.
\end{solution}

\begin{solution}{exo:b3:field-lab-methods-3:8}
(a)~$F = 80/8 = 10$ e $\phi = F^{-1/m} = 10^{-1/2} = \num{0.316}$.
(b)~$10^{-1/1.8} = \num{0.278}$ e $10^{-1/2.2} = \num{0.351}$. Uma variação
de \num{0.2} em $m$ desloca a porosidade em cerca de quatro pontos percentuais;
as resistividades, medidas com precisão de alguns por cento, a deslocam muito menos. O
expoente de cimentação controla a incerteza, e ele deve ser calibrado
com amostras da mesma rocha.
\end{solution}

\begin{solution}{exo:b3:field-lab-methods-3:9}
(a)~Para $s = 2h_1$, $(2nh_1/s)^2 = n^2$, logo os termos são $k^n(1 +
n^2)^{-3/2}$: $0.3536k$, $0.0894k^2$, $0.0316k^3$. Para $k = 0.5$:
$\num{0.1768}$, $\num{0.0224}$, $\num{0.0040}$, e $\rho_a/\rho_1 \approx
1 + 2\times0.2031 = \num{1.406}$ (a soma completa dá \num{1.409}). Para $k
= -0.5$ os termos alternam de sinal, $-0.1768 + 0.0224 - 0.0040$, e
$\rho_a/\rho_1 \approx \num{0.683}$ (soma completa \num{0.685}). (b)~Para $s \ll
h_1$ todos os termos se anulam e $\rho_a \to \rho_1$; para $s \gg h_1$ cada
colchete tende a~1 e $1 + 2\sum k^n = 1 + 2k/(1 - k) = (1 + k)/(1 -
k) = \rho_2/\rho_1$. (c)~$\rho_2/\rho_1 = (1 + k)/(1 - k) = 3$ e $1/3$.
\end{solution}

\begin{solution}{exo:b3:field-lab-methods-3:10}
(a)~Na ordem AMNB, com espaçamento $a$: $AM = a$, $BM = 2a$, $AN = 2a$,
$BN = a$, logo o colchete vale $2/a - 1/a = 1/a$ e $K = 2\pi a$.
(b)~Com A e B em $\mp s$ e M e N em $\mp l$: $AM = BN = s - l$,
$BM = AN = s + l$, o colchete vale $2/(s - l) - 2/(s + l) = 4l/(s^2 -
l^2)$, e $K = \pi(s^2 - l^2)/(2l)$. Quando $l \ll s$, $\Delta V_{MN}
\approx 2l\,E$ com $E = 2(-\mathrm{d}V/\mathrm{d}r)$ no centro, o
campo dos dois eletrodos de corrente; então $\rho_a = K\Delta V/I \to
\pi s^2/(2l)\times 4l(-\mathrm{d}V/\mathrm{d}r)/I = 2\pi s^2(-
\mathrm{d}V/\mathrm{d}r)/I$. (c)~A razão entre o fator aproximado e o
exato é $s^2/(s^2 - l^2) = 1/(1 - 0.01) = 1.010$: um por cento.
\end{solution}

\begin{solution}{exo:b3:field-lab-methods-3:11}
(a)~Leitura: \qty{0.010}{\milli\gal}. Altura:
$0.1967\times0.05 = \qty{0.010}{\milli\gal}$. Densidade: a correção de Bouguer
é $0.04191\,h\,d$ com $d$ em \unit{g.cm^{-3}}, logo
$0.04191\times50\times0.1 = \qty{0.210}{\milli\gal}$. Terreno:
\qty{0.200}{\milli\gal}. Total: $\sqrt{0.010^2 + 0.010^2 + 0.210^2 +
0.200^2} = \qty{0.29}{\milli\gal}$. (b)~A densidade e o terreno
dominam, cada um cerca de vinte vezes o erro de leitura. Para detectar
\qty{0.1}{\milli\gal}, a densidade da laje deve ser medida em
amostras ou ajustada aos dados (o valor que menos reflete a
topografia), e o terreno deve ser corrigido a partir de um modelo de elevação de alta resolução.
(c)~$\sqrt{0.010^2 + 0.010^2 + 0.210^2} = \qty{0.21}{\milli\gal}$.
\end{solution}

\begin{solution}{exo:b3:field-lab-methods-3:12}
(a)~A massa que se ajusta à leitura central é $\Delta m = \Delta g(0)z^2/G
= \num{-4.16e-5}z^2/\num{6.674e-11}$: \qty{-5.61e10}{kg} a
\qty{300}{m}, \qty{-9.97e10}{kg} a \qty{400}{m}, \qty{-1.56e11}{kg} a
\qty{500}{m}. As leituras previstas em $x = \pm300$ e $\pm600$ são
\num{-1.47} e \num{-0.37}, \num{-2.13} e \num{-0.71}, \num{-2.62}
e \qty{-1.09}{\milli\gal}. Os \omterm{def:b3:field-lab-methods-3:inversion-practice}{desajustes} são $\chi^2 = 1238$, $3.7$ e
$856$. (b)~Apenas o modelo a \qty{400}{m}, com $\chi^2 = 3.7$ para cinco
dados, ajusta-se dentro dos erros. (c)~Sim: qualquer corpo esfericamente simétrico
com a mesma massa e centro produz exatamente a mesma anomalia, e outras
formas podem caber dentro dos erros. Uma densidade medida em amostras, em um
furo de sondagem ou outro método decidiria.
\end{solution}

\begin{solution}{pb:b3:field-lab-methods-3:1}
\textbf{1.} $(2846.708 - 2846.512)/\qty{2.0}{h} = \qty{0.098}{\milli\gal.h^{-1}}$.

\textbf{2.} $2851.075 - 0.098\times0.5 - 2846.512 = \qty{4.514}{\milli\gal}$;
$2850.497 - 0.088 - 2846.512 = \qty{3.897}{\milli\gal}$;
$2849.392 - 0.127 - 2846.512 = \qty{2.753}{\milli\gal}$.

\textbf{3.} As alturas acima da base são $-27.1$, $-30.6$ e
\qty{-18.1}{m}; somando $0.1967\,\Delta h$ obtêm-se as anomalias de Bouguer
$\num{-0.817}$, $\num{-2.122}$ e \qty{-0.808}{\milli\gal}. As
estações estão abaixo da base e mais próximas do centro da Terra,
por isso suas leituras brutas são maiores; a correção elimina isso e
deixa um mínimo sobre o granito.

\textbf{4.} A parcela extra $0.098 - 0.040 = \qty{0.058}{\milli\gal.h^{-1}}$,
\qty{0.12}{\milli\gal} ao longo do circuito, é a maré terrestre sólida. Em duas
horas ela varia quase linearmente, de modo que uma reta que passa pelas duas
leituras na base a remove junto com a deriva.

\textbf{5.} Os erros são $\num{0.008}$, $\num{0.009}$ e
\qty{0.017}{\milli\gal}. Arredondar cada altura para \qty{0.1}{m} pode tornar
uma diferença de altura errada em até \qty{0.1}{m}, ou seja,
\qty{0.020}{\milli\gal}: o arredondamento explica os erros; a maré
residual após a reta, sendo inferior a \qty{0.001}{\milli\gal}.

\textbf{6.} $0.01/0.1967 = \qty{5.1}{cm}$. Os \qty{3.3}{cm} dos manuais correspondem a
$0.01/0.3086$: a anomalia ar livre não depende da densidade
escolhida para a correção de Bouguer, e a altura também deve servir a ela.

\textbf{7.} $x_{1/2} = (3.952 - 3.048)/2 = \qty{0.452}{km}$, e $z = 1.305
\times452 = \qty{590}{m}$ abaixo do solo.

\textbf{8.} $\Delta m = \Delta g\,z^2/G = \num{-2.293e-5}\times590^2/
\num{6.674e-11} = \qty{-1.20e11}{kg}$. Com $\Delta\rho =
\qty{-130}{kg.m^{-3}}$, seu volume é \qty{9.2e8}{m^3} e seu raio
\qty{603}{m}, maior que a profundidade de seu centro: a esfera
ultrapassaria o solo. Nenhuma esfera com o contraste de densidade do granito
se ajusta; o corpo deve chegar perto da superfície.

\textbf{9.} O granito se estende da superfície a grande profundidade. Sua parte rasa
próxima às estações, torna o pico estreito, e a regra, que
conhece apenas a largura, coloca uma única esfera a uma profundidade moderada; a
parte profunda acrescenta flancos largos que a esfera não tem (a figura do
perfil), e a massa que eles contêm não é detectada.

\textbf{10.} $2\pi G\Delta\rho = 2\pi\times\num{6.674e-11}\times(-130) =
\qty{-5.45e-8}{s^{-2}}$. Com $z_1 = 0$, $z_2 = \qty{1645}{m}$, $a =
\qty{461}{m}$: o colchete é $1645 + 461 - 1708.4 = \qty{397.6}{m}$,
e $\Delta g = \qty{-2.17}{\milli\gal}$. Para uma profundidade ilimitada o
colchete tende a $a$, e $\Delta g = \qty{-2.51}{\milli\gal}$.

\textbf{11.} O colchete deve ser $2.29/\num{5.45e-3} = \qty{420}{m}$, o que um
piso a cerca de \qty{2.6}{km} fornece. Qualquer piso entre cerca de
\qty{1.6}{km} e uma profundidade ilimitada altera o pico em apenas algumas
décimas de miligal, e um corpo ligeiramente mais largo ou mais denso dá o mesmo
pico com um piso mais raso: o piso não é resolvido.

\textbf{12.} Com $z_1 = \qty{6}{m}$: $1639 + 461.0 - 1708.4 = 391.6$, e
$\Delta g = \qty{-2.14}{\milli\gal}$, \qty{0.03}{\milli\gal} a menos que
com o topo no solo. Isso está abaixo dos erros de terreno e de densidade
do levantamento: a gravidade não consegue dizer se o granito atinge
a superfície sob o aluvião. A sondagem consegue.

\textbf{13.} $K = \pi(20^2 - 1^2)/(2\times1) = \qty{626.7}{m}$, e $\Delta V =
\rho_a I/K = 821\times0.050/626.7 = \qty{65.5}{mV}$.

\textbf{14.} $V = 2\times1/\sqrt{17.1^2 - 10.7^2} = 2/13.34 =
\qty{0.150}{m.ns^{-1}}$, $\varepsilon_r = (0.2998/0.150)^2 = 4.0$, uma
areia seca; o tubo está a $0.150\times10.7/2 = \qty{0.80}{m}$ de profundidade.

\textbf{15.} $0.150\times20.0/2 = \qty{1.50}{m}$.

\textbf{16.} $20/821 = \num{0.0244}$, $30/1309 = \num{0.0229}$, $40/1684 =
\num{0.0238}$: $S \approx \qty{0.0237}{S}$.

\textbf{17.} A camada seca contribui com $1.5/600 = \qty{0.0025}{S}$, a
saturada $\qty{0.0212}{S}$, logo sua espessura é $0.0212\times222 =
\qty{4.7}{m}$, e o granito fica $1.5 + 4.7 = \qty{6.2}{m}$ abaixo
da superfície. A estimativa é alta: o granito não é infinitamente
resistivo, de modo que, nestes espaçamentos, $\rho_a$ ainda fica abaixo da reta
$s/S$, e $s/\rho_a$ superestima $S$. (O modelo da região tem
\qty{6}{m} de aluvião.)

\textbf{18.} $\phi = (20/222)^{1/2} = \num{0.30}$. Com $m = 1.3$, $\phi =
0.0901^{1/1.3} = \num{0.16}$; com $m = 2.8$, $\phi = \num{0.42}$.

\textbf{19.} Subtraindo a base da linha obtêm-se, de oeste para leste:
3.4, 4.1, 4.8, 5.9, 7.4, 9.2, 11.5, 14.2, 17.2, 19.8, 21.0, 20.1, 17.6,
14.5, 11.5, 9.1, 7.2, 5.7, 4.6, 3.8 e \qty{3.1}{nT}.

\textbf{20.} Sem correção, o pico é \qty{28.0}{nT} em \qty{2.20}{km}, e
a metade dele, \qty{14.0}{nT}, nunca é atingida a leste: a linha termina
em \qty{15.2}{nT}, porque o campo subiu \qty{12}{nT} durante a
caminhada.

\textbf{21.} O pico é \qty{21.0}{nT} em \qty{2.20}{km}; a metade dele é
\qty{10.5}{nT}. A oeste: entre \num{9.2} em \num{1.95} e \num{11.5} em
\qty{2.00}{km}, em $1.95 + 0.05\times1.3/2.3 = \qty{1.978}{km}$. A leste:
entre \num{11.5} em \num{2.40} e \num{9.1} em \qty{2.45}{km}, em
$2.40 + 0.05\times1.0/2.4 = \qty{2.421}{km}$. A meia-largura é $(2.421
- 1.978)/2 = \qty{0.221}{km}$.

\textbf{22.} Uma anomalia simétrica em um campo com inclinação de \ang{65} significa que apenas
a magnetização vertical atua: o corpo tem direção paralela ao
meridiano magnético, norte--sul. A regra
da \cref{prop:b3:field-lab-methods-3:dyke} supõe um dique fino, vertical,
bidimensional que se estende a grande profundidade, magnetizado por indução, cujo
topo atua como uma linha de carga magnética.

\textbf{23.} $\kappa t = 2\pi z\,\Delta T_\mathrm{max}/(F\sin^2 I) = 2\pi\times
221\times21.0/(48000\times0.8214) = \qty{0.74}{m}$. Com $\kappa =
0.03$, $t = \qty{25}{m}$. Apenas o produto é fixado: um dique duas vezes
mais espesso com metade da suscetibilidade dá a mesma anomalia.

\textbf{24.} A anomalia tem \qty{443}{m} de largura à meia altura do pico e as estações
estão separadas por \qty{50}{m}: nove leituras atravessando-a, três vezes o
mínimo da regra.

\textbf{25.} O topo do dique está
\textbf{\qty[text-series-to-math=true]{220}{m}} abaixo do solo. A
modelo geológico situa a discordância a \qty{222}{m} de profundidade neste ponto:
o dique sobe até a base da cobertura e é truncado ali, mais antigo
que o arenito triássico que o soterra.
\end{solution}
